\originalsection{第12次习题课}
\sourceinfo{MIT 18.04, Spring 2018, Recitation 12; 题面 \nolinkurl{mit18_04s18_recit12-handout.pdf}, PDF第1页; 解答 \nolinkurl{mit18_04s18_recit12-solutions.pdf}, PDF第1页. 课程作者 Jeremy Orloff, 页内署名 Vishesh Jain, 2018, CC BY-NC-SA 4.0}
\begin{exercise}
\textbf{原题 Rec12.1。}\label{ass:rec12-1}
用 Rouch\'e 定理证明 $z^5+3z+1$ 的全部5个零点都位于圆 $C_2=\{z:|z|=2\}$ 内, 按重数计.
\end{exercise}
\begin{solution}
编者补充 (AI): 官方解答仅指向讲义例11.7. 在 $|z|=2$ 上, $|z^5|=32$, $|3z+1|\le7<32$. Rouch\'e 定理使 $z^5+3z+1$ 与 $z^5$ 在 $|z|<2$ 中有相同零点数, 即5个, 按重数计. 不等式也说明圆上没有零点; 五次多项式至多有5个根, 所以所有根均已计入.
\end{solution}
\begin{exercise}
\textbf{原题 Rec12.2。}\label{ass:rec12-2}
用 Rouch\'e 定理证明 $z+3+2\ee^z$ 在左半平面恰有一个零点.
\end{exercise}
\begin{solution}
编者补充 (AI): 官方解答仅指向讲义例11.8. 取由左半圆弧 $|z|=R,\Re z\le0$ 与虚轴线段构成的正向闭轮廓, $R>5$. 在虚轴上 $|z+3|\ge3>2=|2\ee^z|$; 在左半圆弧上 $|z+3|\ge R-3>2\ge|2\ee^z|$. 因而 Rouch\'e 定理说明半圆内 $z+3+2\ee^z$ 与 $z+3$ 零点数相同, 为1. 对任何 $\Re z\le0$ 的零点, 又有 $|z+3|=2\ee^{\Re z}\le2$, 故 $|z|\le5$. 因而较大半圆不会遗漏左半平面的零点. 虚轴上的严格不等式排除边界零点, 结论恰为左开半平面一个零点, 按重数计, 从而该零点简单.
\end{solution}
\begin{exercise}
\textbf{原题 Rec12.3。}\label{ass:rec12-3}
用 Rouch\'e 定理再证明代数基本定理: 首一多项式 $P(z)=z^n+a_{n-1}z^{n-1}+\cdots+a_0$ 在复平面恰有 $n$ 个根, 按重数计.
\end{exercise}
\begin{solution}
编者补充 (AI): 官方解答仅指向例11.8之后的定理. 若 $n=0$, 非零常数没有根; 以下 $n\ge1$. 选 $R>1$ 且 $R>\sum_{j=0}^{n-1}|a_j|$. 在 $|z|=R$ 上,
\[
|P(z)-z^n|\le\sum_{j=0}^{n-1}|a_j|R^j
\le R^{n-1}\sum_{j=0}^{n-1}|a_j|<R^n=|z^n|.
\]
Rouch\'e 定理给 $P$ 在圆内有 $n$ 个零点, 按重数计. 对任意 $|z|\ge R$, 相同估计得到 $|P(z)-z^n|<|z^n|$, 所以圆外也不可能有根. 因而复平面总根数恰为 $n$, 不需要预先使用代数基本定理.
\end{solution}
\begin{exercise}
\textbf{原题 Rec12.4。}\label{ass:rec12-4}
设 $G$ 亚纯, $H=G/(1+G)$. 对闭曲线 $\gamma$, 假设 $G\circ\gamma$ 不经过 $-1$, $G$ 在 $\gamma$ 上没有极点, 证明 $P_{H,\gamma}=P_{G,\gamma}+\Ind(G\circ\gamma,-1)$, 其中 $P$ 表示所围区域内极点数, 按重数计.
\end{exercise}
\begin{solution}
据官方解答翻译补全; 曲线条件与极点处的消去说明为编者补充 (AI). 若 $P$ 表示普通内部计数, 须取 $\gamma$ 为正向分段光滑 Jordan 曲线, 且 $G$ 在其闭内部的某邻域亚纯. 原题只写“闭曲线”; 对一般闭曲线, 以下 $Z,P$ 要改为按曲线绕数加权的计数, 才保持公式成立.

在 $G$ 的极点 $a$ 处, $1/G$ 解析且为零, 而 $H=1/(1+1/G)$ 可解析延伸, 值为1. 因而 $G$ 的极点不会成为 $H$ 的极点. 在 $1+G$ 的零点处 $G=-1\ne0$, 所以 $H$ 极点阶数正好等于该零点阶数, 得 $P_H=Z_{1+G}$. 另一方面, 加上常数1不改变 $G$ 的极点阶数, 所以 $P_{1+G}=P_G$. 辐角原理给
\[
Z_{1+G}-P_G=\Ind((1+G)\circ\gamma,0)=\Ind(G\circ\gamma,-1).
\]
代入 $P_H=Z_{1+G}$, 就得到所求关系. $H$ 在原式分母无穷大的点采用上述亚纯延伸, 不能把这些点机械计为极点.
\end{solution}
\begin{exercise}
\textbf{原题 Rec12.5。}\label{ass:rec12-5}
在原点有一个源, 两堵墙为 $y=1$ 和 $y=-1$. 用镜像法需要多少个像源? 所得复势是什么?
\end{exercise}
\begin{solution}
据官方解答翻译补全; 原级数的纠错、正则化和墙上速度核验为编者补充 (AI). 先在 $(0,2)$ 放置同强度源, 使上墙成为流线. 为恢复下墙条件须加入 $(0,-2),(0,-4)$; 再反射到上墙得到 $(0,4),(0,6)$, 如此继续. 两次墙反射的复合是竖直平移4, 因而闭合的源集合为 $z=2n\ii$, $n\in\Z$. 除真实源 $n=0$ 外, 需要无穷多个同强度像源.

官方解答印出的势为
\[
\Phi_{\rm original}(z)=\log z+\sum_{n\ge1}\log(z+2n\ii)+\sum_{n\le-1}\log(z-2n\ii).
\]
\textbf{第二个求和的 $n\le-1$ 使它与第一个求和重复了负侧源, 丢失正侧源; 即使修正源位置, 未正则化的对数级数也不收敛.} 正确的形式源项应来自 $\log(z-2n\ii)+\log(z+2n\ii)$, $n\ge1$. 两项配对后去掉与 $z$ 无关的常数, 可定义局部复势
\[
\Phi(z)=\Log z+\sum_{n=1}^\infty\Log\left(1+\frac{z^2}{4n^2}\right)+C.
\]
这里在不含源点的单连通小域内选择相容对数分支; 势全局可为多值, 速度是单值. 任意紧集上, 大 $n$ 的对数项为 $O(n^{-2})$, 故尾级数局部一致绝对收敛. 逐项求导得
\[
\Phi'(z)=\frac1z+\sum_{n=1}^\infty\frac{2z}{z^2+4n^2}
=\frac1z+\sum_{n=1}^\infty\left(\frac1{z-2n\ii}+\frac1{z+2n\ii}\right).
\]
每个像源留数为1, 正是所需的同强度源; 配对顺序也说明不能任意拆开两个无穷和.

为识别闭式, 利用第13章由 Gamma 反射公式与 Euler 乘积推出的正弦乘积 $\sin(\pi\zeta)=\pi\zeta\prod_{n\ge1}(1-\zeta^2/n^2)$, 代入 $\zeta=z/(2\ii)$, 得
\[
z\prod_{n=1}^\infty\left(1+\frac{z^2}{4n^2}\right)=\frac2\pi\sinh\frac{\pi z}2.
\]
因此局部势也可写为 $\Log((2/\pi)\sinh(\pi z/2))+C$, 等价于 $\Log\sinh(\pi z/2)$ 加常数, 且 $\Phi'=(\pi/2)\coth(\pi z/2)$. 在 $z=x\pm\ii$ 上, $\coth(\pi x/2\pm\pi\ii/2)=\tanh(\pi x/2)\in\R$. 流速 $\mathbf F=(\Re\Phi',-\Im\Phi')$ 的法向分量 $F_y=0$, 两墙确为流线. 常数及对数分支不会影响此结论.
\begin{center}
\begin{tikzpicture}[scale=.6,>=Stealth]
\draw[->](-3,0)--(3,0)node[right]{$x$};\draw[->](0,-6.7)--(0,6.7)node[above]{$y$};
\draw[thick](-2.8,1)--(2.8,1)node[right]{墙 $y=1$};\draw[thick](-2.8,-1)--(2.8,-1)node[right]{墙 $y=-1$};
\fill(0,0)circle(2pt)node[below right]{真实源};
\foreach \y in {-6,-4,-2,2,4,6}{\draw[fill=white](0,\y)circle(2pt) node[right]{$\y\ii$};}
\node at(0,6.35){$\vdots$};\node at(0,-6.35){$\vdots$};
\node[align=left] at(-2.2,3.8){同强度\\像源};
\end{tikzpicture}
\end{center}

\end{solution}
